%% =========================================================
%% Example: group matrix product operators from a normalized
%% three-cocycle on periodic chains, with the instances for the
%% group of order two (the decorated CZX operator), the cyclic
%% groups with their standard cocycles, and Z2 x Z2 with the
%% mixed cocycle of type II.
%% =========================================================

\section{Group matrix product operators from a three-cocycle}%
\label{sec:asymex_group_cocycle_mpo}

Garre-Rubio, Lootens and Molnár~\cite{GarreRubio2023MPOSym} build, from a
finite group $G$ and a three-cocycle $\omega$, a matrix product operator
representation of $G$ on a periodic chain whose sites carry the regular
representation $\C[G]$. Each operator is a left shift by $g$ on every site,
decorated by a diagonal two-site phase gate determined by $\omega$. For the
cyclic groups, Garre-Rubio and Schuch~\cite{GarreRubioSchuch2024DomainWall}
give one explicit cocycle in each cohomology class.
The periodic operator below places one left shift per site following the
local tensor of~\cite{GarreRubio2023MPOSym}, rather than the doubled shift of
its periodic display~\cite{gap:glm23_pbc_group_mpo_single_shift}.

\begin{definition}[The tensor of a three-cochain]
    \label{def:asymex_group_cocycle_tensor}
    \lean{MPOTensor.GroupCocycle.wGate,
        MPOTensor.GroupCocycle.leftShift,
        MPOTensor.GroupCocycle.leftShift_apply,
        MPOTensor.GroupCocycle.siteShift,
        MPOTensor.GroupCocycle.siteShift_apply,
        MPOTensor.GroupCocycle.tensor,
        MPOTensor.GroupCocycle.tensor_apply,
        MPOTensor.GroupCocycle.family,
        MPOTensor.GroupCocycle.shift,
        MPOTensor.GroupCocycle.shift_apply,
        MPOTensor.GroupCocycle.phase}
    \leanok
    \uses{def:mpo_tensor, def:mpo_family, def:mpug_scalar_cochains}
    Let $G$ be a finite group and let
    $\omega\colon G^3\to\C^\times$ be a scalar three-cochain. On
    $\C[G]\otimes\C[G]$ define the diagonal gate and the left shift by
    \begin{align}
        W_g\ket{k,l} & =\omega(g,l,l^{-1}k)\ket{k,l}, &
        L_g\ket{k}   & =\ket{gk}.
        \notag
    \end{align}
    The tensor $\hat T_g$ has physical and bond indices in $G$, with
    \begin{align}
        (\hat T_g^{\,a,b})_{l,r}
         & =\delta_{a,gb}\,\delta_{r,b}\,\omega(g,b,b^{-1}l),
        \label{eq:asymex_group_cocycle_tensor}
    \end{align}
    where $a$ is the output, $b$ the input, and $l,r$ the incoming and
    outgoing bonds. For a configuration $t=(t_0,\ldots,t_{N-1})\in G^N$
    on a periodic chain of $N>0$ sites, with $t_N=t_0$, write
    $gt=(gt_0,\ldots,gt_{N-1})$ and
    \begin{align}
        \phi_g(t)
         & =\prod_{i=0}^{N-1}\omega(g,t_{i+1},t_{i+1}^{-1}t_i).
        \label{eq:asymex_group_cocycle_phase}
    \end{align}
    The group elements are identified with the basis labels
    $\{0,\ldots,|G|-1\}$ by a fixed bijection, and $U_g=O_N(\hat T_g)$
    denotes the periodic operator.
\end{definition}

\begin{theorem}[The periodic operator is monomial]
    \label{thm:asymex_group_cocycle_kernel}
    \lean{MPOTensor.GroupCocycle.mpo_tensor_apply,
        MPOTensor.GroupCocycle.mpo_tensor,
        MPOTensor.GroupCocycle.mpo_tensor_apply_eq_leftShift_wGate}
    \leanok
    \uses{def:asymex_group_cocycle_tensor}
    For every scalar three-cochain $\omega$, every $g\in G$, every $N>0$,
    and all $s,t\in G^N$,
    \begin{align}
        U_g(s,t) & =\delta_{s,gt}\,\phi_g(t).
        \notag
    \end{align}
    Equivalently,
    $U_g=\bigl(\bigotimes_iL_g\bigr)\prod_iW_g^{i,i+1}$, the gates acting on
    the neighboring pairs $(t_i,t_{i+1})$ before the shifts.
\end{theorem}

\begin{proof}
    \leanok
    Expand the trace over cyclic bond configurations $b$. The factor
    $\delta_{r,b}$ in~(\ref{eq:asymex_group_cocycle_tensor}) forces
    $b_{i+1}=t_i$, so only $b_i=t_{i-1}$ contributes. Its contribution
    vanishes unless $s_i=gt_i$ at every site, and otherwise equals
    $\prod_i\omega(g,t_i,t_i^{-1}t_{i-1})$, which is $\phi_g(t)$ after the
    cyclic reindexing $i\mapsto i+1$. The gates are diagonal, so the entry
    of $\bigl(\bigotimes_iL_g\bigr)\prod_iW_g^{i,i+1}$ at $(s,t)$ is the
    same product.
\end{proof}

\begin{theorem}[Product law]
    \label{thm:asymex_group_cocycle_mul}
    \lean{MPOTensor.GroupCocycle.mpo_tensor_mul,
        MPOTensor.GroupCocycle.phase_mul,
        MPOTensor.GroupCocycle.shift_mul}
    \leanok
    \uses{thm:asymex_group_cocycle_kernel, def:mpug_three_cocycle}
    If $\omega$ is a three-cocycle, then $U_gU_h=U_{gh}$ for all
    $g,h\in G$ and every $N>0$.
\end{theorem}

\begin{proof}
    \leanok
    Both sides are monomial with the shift by $gh$, so it suffices that
    $\phi_g(ht)\phi_h(t)=\phi_{gh}(t)$. The cocycle
    equation~(\ref{eq:mpug_three_cocycle}) at $(g,h,a,a^{-1}x)$ gives, at
    each site,
    \begin{align}
        \omega(g,ha,a^{-1}x)\,\omega(h,a,a^{-1}x)
         & =\omega(gh,a,a^{-1}x)\,
        \frac{\omega(g,h,x)}{\omega(g,h,a)}.
        \notag
    \end{align}
    With $a=t_{i+1}$ and $x=t_i$ the correction factors telescope around
    the ring.
\end{proof}

\begin{theorem}[Unit law]
    \label{thm:asymex_group_cocycle_one}
    \lean{MPOTensor.GroupCocycle.mpo_tensor_one,
        MPOTensor.GroupCocycle.wGate_one,
        MPOTensor.GroupCocycle.shift_one}
    \leanok
    \uses{thm:asymex_group_cocycle_kernel, def:mpug_scalar_cochains}
    If $\omega$ is normalized, then $W_e=1$ and $U_e=1$ for every $N>0$.
\end{theorem}

\begin{proof}
    \leanok
    Normalization in the first argument makes every factor of $W_e$ and of
    $\phi_e(t)$ equal to one, and the shift by $e$ is the identity.
\end{proof}

\begin{theorem}[Unitarity and the adjoint law]
    \label{thm:asymex_group_cocycle_unitary}
    \lean{MPOTensor.GroupCocycle.mpo_tensor_mem_unitaryGroup,
        MPOTensor.GroupCocycle.star_phase_mul_phase,
        MPOTensor.GroupCocycle.isMPUPos,
        MPOTensor.GroupCocycle.mpo_tensor_conjTranspose,
        MPOTensor.GroupCocycle.family_operator_laws}
    \leanok
    \uses{thm:asymex_group_cocycle_kernel, thm:asymex_group_cocycle_mul,
        thm:asymex_group_cocycle_one, def:mpug_positive_mpu}
    If every value of $\omega$ has modulus one, then $U_g$ is unitary for
    every $g\in G$ and every $N>0$, so $\hat T_g$ is a matrix product
    unitary on nonempty chains. If moreover $\omega$ is a normalized
    three-cocycle, then $U_g^\dagger=U_{g^{-1}}$ for every $N>0$.
\end{theorem}

\begin{proof}
    \leanok
    A monomial matrix whose nonzero entries have modulus one is unitary.
    For the adjoint law, $U_g^\dagger=U_g^\dagger U_gU_{g^{-1}}=U_{g^{-1}}$
    by Theorems~\ref{thm:asymex_group_cocycle_mul}
    and~\ref{thm:asymex_group_cocycle_one}.
\end{proof}

The single-site tensors are not injective: their $|G|^2$ physical matrices
span a space of dimension at most $|G|$. The operator laws above are
therefore stated directly rather than through an injective representation.

\subsection{The group of order two and the decorated CZX operator}

For $G=\Z_2=\{e,s\}$ with the nontrivial cocycle $\omega_1$ of
Definition~\ref{def:asymex_z2_cocycles}, the source identifies the gate as
$CZ(1\otimes Z)$, the shift as $X$, and the operator as
$\prod_iCZ_{i,i+1}Z_i\prod_iX_i$~\cite{GarreRubio2023MPOSym}. Identify
$e,s$ with the qubit labels $0,1$ and write $\bar t$ for the bitwise
complement of $t\in\{0,1\}^N$ and $q(s)=\sum_is_is_{i+1}$.

\begin{theorem}[The decorated CZX operator]
    \label{thm:asymex_group_cocycle_z2}
    \lean{MPOTensor.GroupCocycle.bitEquiv,
        MPOTensor.GroupCocycle.czxDecorated,
        MPOTensor.GroupCocycle.wGate_cyclicTwo,
        MPOTensor.GroupCocycle.mpo_tensor_cyclicTwo}
    \leanok
    \uses{def:asymex_z2_cocycles, thm:asymex_group_cocycle_kernel}
    For $\omega=\omega_1$ and $g=s$, the gate is
    $W_s\ket{k,l}=(-1)^{kl+l}\ket{k,l}$, which is $CZ(1\otimes Z)$. For
    every $N>0$,
    \begin{align}
        U_s(u,t) & =\delta_{u,\bar t}\,(-1)^{q(u)+\sum_iu_i},
        \notag
    \end{align}
    so that $U_s=\prod_iCZ_{i,i+1}Z_i\prod_iX_i$.
\end{theorem}

\begin{proof}
    \leanok
    On bits, $\omega_1(s,l,l^{-1}k)=(-1)^{kl+l}$, checked on the four
    pairs $(k,l)$. By Theorem~\ref{thm:asymex_group_cocycle_kernel} the
    kernel is supported on $u=\bar t$ with phase
    $\prod_i(-1)^{t_it_{i+1}+t_{i+1}}$. On complemented bits,
    $(-1)^{kl+l}=(-1)^{\bar k\bar l+\bar k}$, and the product becomes
    $(-1)^{q(u)+\sum_iu_i}$ after reindexing the second factor.
\end{proof}

\subsection{The cyclic groups}

For $n\geq1$ and $j\in\N$, Garre-Rubio and
Schuch~\cite[Section~IV.D]{GarreRubioSchuch2024DomainWall} write the
three-cocycles of $\Z_n$ as $\omega_j$ below, with $j=0,\ldots,n-1$
labelling the cohomology classes.

\begin{definition}[The cocycles of the cyclic groups]
    \label{def:asymex_group_cocycle_cyclic}
    \lean{TNLean.Algebra.ScalarThreeCochain.cyclicCocycle,
        TNLean.Algebra.ScalarThreeCochain.rootOfUnity,
        TNLean.Algebra.ScalarThreeCochain.rootOfUnity_pow,
        MPOTensor.GroupCocycle.residueEquiv}
    \leanok
    \uses{def:mpug_scalar_cochains}
    Represent residues of $\Z_n$ by $\{0,\ldots,n-1\}$ and write
    $[b+c]$ for the representative of $b+c$. Define
    \begin{align}
        \omega_j(a,b,c)
         & =\exp\Bigl\{\frac{2\pi\mathrm{i}\,j}{n^2}\,a\bigl(b+c-[b+c]\bigr)\Bigr\},
        \notag
    \end{align}
    and let $\zeta=e^{2\pi\mathrm{i}/n}$.
\end{definition}

\begin{theorem}[The cyclic instances]
    \label{thm:asymex_group_cocycle_cyclic}
    \lean{TNLean.Algebra.ScalarThreeCochain.cyclicCocycle_val,
        TNLean.Algebra.ScalarThreeCochain.cyclicCocycle_isCocycle,
        TNLean.Algebra.ScalarThreeCochain.cyclicCocycle_isNormalized,
        TNLean.Algebra.ScalarThreeCochain.cyclicCocycle_norm,
        MPOTensor.GroupCocycle.mpo_tensor_cyclic_apply,
        MPOTensor.GroupCocycle.cyclic_operator_laws}
    \leanok
    \uses{def:asymex_group_cocycle_cyclic, def:mpug_three_cocycle,
        thm:asymex_group_cocycle_kernel, thm:asymex_group_cocycle_mul,
        thm:asymex_group_cocycle_one, thm:asymex_group_cocycle_unitary}
    For every $n\geq1$ and $j\in\N$, the value is
    $\omega_j(a,b,c)=\zeta^{ja\lfloor(b+c)/n\rfloor}$, and $\omega_j$ is a
    normalized three-cocycle of $\Z_n$ with values of modulus one. For every
    $g\in\Z_n$, every $N>0$ and all $s,t\in\Z_n^N$,
    \begin{align}
        U_g(s,t)
         & =\delta_{s,g+t}\prod_{i=0}^{N-1}
        \zeta^{jg\lfloor(t_{i+1}+[t_i-t_{i+1}])/n\rfloor}.
        \notag
    \end{align}
    Every $\hat T_g$ is a matrix product unitary on nonempty chains, and
    $U_0=1$, $U_gU_h=U_{g+h}$ and $U_g^\dagger=U_{-g}$ for every $N>0$.
\end{theorem}

\begin{proof}
    \leanok
    Since $b+c-[b+c]=n\lfloor(b+c)/n\rfloor$, the exponential is the stated
    power of $\zeta$. The cocycle equation reduces modulo $n$ to the carry
    identity
    \begin{align}
        \Bigl\lfloor\frac{k+l}{n}\Bigr\rfloor
        +\Bigl\lfloor\frac{h+[k+l]}{n}\Bigr\rfloor
         & =\Bigl\lfloor\frac{h+k}{n}\Bigr\rfloor
        +\Bigl\lfloor\frac{[h+k]+l}{n}\Bigr\rfloor,
        \notag
    \end{align}
    both sides counting the multiples of $n$ in $h+k+l$. An identity
    argument makes the exponent zero, and every power of $\zeta$ has
    modulus one. The kernel is
    Theorem~\ref{thm:asymex_group_cocycle_kernel} evaluated on these
    values, and the operator laws are
    Theorems~\ref{thm:asymex_group_cocycle_mul},
    \ref{thm:asymex_group_cocycle_one}
    and~\ref{thm:asymex_group_cocycle_unitary}.
\end{proof}

\subsection{The group \texorpdfstring{$\Z_2\times\Z_2$}{Z2 x Z2}}

For $\Z_2\times\Z_2$ the construction is applied to the mixed three-cocycle
of type II whose exponent couples the first factor to the second, the
cocycle targeted in Section~\ref{sec:asymex_z2z2}. It is the type-II
component in the decomposition of the three-cocycles of $\Z_2\times\Z_2$
printed in~\cite[Section~6.1]{GarreRubio2023MPOSym}, where it takes the
value $-1$ exactly on the triples $(ab^i,a^jb,a^kb)$ with $a=(1,0)$ and
$b=(0,1)$. That source does not apply the periodic construction to it.

\begin{definition}[The mixed cocycle of type II]
    \label{def:asymex_group_cocycle_klein}
    \lean{TNLean.Algebra.ScalarThreeCochain.kleinCocycle,
        MPOTensor.GroupCocycle.kleinEquiv}
    \leanok
    \uses{def:mpug_scalar_cochains}
    For $a=(a_1,a_2)$, $b=(b_1,b_2)$ and $c=(c_1,c_2)$ in $\Z_2\times\Z_2$,
    define $\omega(a,b,c)=(-1)^{a_1b_2c_2}$.
\end{definition}

\begin{theorem}[The instance for the group of order four]
    \label{thm:asymex_group_cocycle_klein}
    \lean{TNLean.Algebra.ScalarThreeCochain.kleinCocycle_isCocycle,
        TNLean.Algebra.ScalarThreeCochain.kleinCocycle_isNormalized,
        TNLean.Algebra.ScalarThreeCochain.kleinCocycle_norm,
        MPOTensor.GroupCocycle.mpo_tensor_klein_apply,
        MPOTensor.GroupCocycle.klein_operator_laws}
    \leanok
    \uses{def:asymex_group_cocycle_klein, def:mpug_three_cocycle,
        thm:asymex_group_cocycle_kernel, thm:asymex_group_cocycle_mul,
        thm:asymex_group_cocycle_one, thm:asymex_group_cocycle_unitary}
    The function $\omega$ is a normalized three-cocycle with values $\pm1$.
    Write $t_i=(x_{i,1},x_{i,2})$. For every $g\in\Z_2\times\Z_2$, every
    $N>0$ and all $s,t$,
    \begin{align}
        U_g(s,t)
         & =\delta_{s,g+t}\,
        (-1)^{\sum_ig_1x_{i+1,2}(x_{i,2}-x_{i+1,2})}.
        \notag
    \end{align}
    Every $\hat T_g$ is a matrix product unitary on nonempty chains, and
    $U_0=1$, $U_gU_h=U_{g+h}$ and $U_g^\dagger=U_{-g}$ for every $N>0$.
\end{theorem}

\begin{proof}
    \leanok
    The exponent $a_1b_2c_2$ is trilinear over $\Z_2$, hence an additive
    three-cocycle, and it vanishes when any argument is zero. The kernel is
    Theorem~\ref{thm:asymex_group_cocycle_kernel} with
    $t_{i+1}^{-1}t_i=t_i-t_{i+1}$, and the operator laws are
    Theorems~\ref{thm:asymex_group_cocycle_mul},
    \ref{thm:asymex_group_cocycle_one}
    and~\ref{thm:asymex_group_cocycle_unitary}.
\end{proof}
